% !PRINTED 424
\weberpage{424}

%<<CONT
and the sign of this depends on $k$. But if we take $k = 1$, and hence $r = r_0$, then all the angles
\[
\frac{2\pi}{n}, \frac{4\pi}{n} \ldots \frac{(n-1)\pi}{n}
\]
lie between zero and $\pi$, and the sines are all positive. We conclude from this that in this case
\weq{14}{P = (-i)^{\frac{n-1}{2}} n^{\frac{n-2}{2}} \sqrt{n}}
and that $\sqrt{n}$ is to be taken positive.

The sign of the product of the $R$ does not depend on the choice of $r$, but only on the nature of $n$. It is therefore of interest to determine by itself the product $Q$, whose sign does depend on the choice of $r$. We obtain it in the following way:

If we multiply the expression (11) by
\[
r \cdot r^2 \cdot r^3 \cdots r^{\frac{n-1}{2}} = r^{\frac{n^2-1}{8}}
\]
and then by
\[
r^{-1} \cdot r^{-2} \cdot r^{-3} \cdots r^{-\frac{n-1}{2}} = r^{-\frac{n^2-1}{8}},
\]
we obtain
\[
r^{\frac{n^2-1}{8}} Q = (r^2 - 1)(r^4 - 1) \cdots (r^{n-1} - 1)
\]
\[
r^{-\frac{n^2-1}{8}} Q = (1 - r^{n-2})(1 - r^{n-4}) \cdots (1 - r).
\]
Now since the exponents $2, 4 \ldots n-1, n-2, n-4 \ldots 1$ together comprise all the numbers $1, 2 \ldots n-1$, multiplication of these two expressions yields:
\[
Q^2 = (-1)^{\frac{n-1}{2}} \Pi(r-1) = (-1)^{\frac{n-1}{2}} n,
\]
hence
\weq{15}{Q = \pm i^{\frac{n-1}{2}} \sqrt{n}.}
Comparison with (13) yields
\weq{16}{2^{\frac{n-1}{2}} \sin\frac{2\pi k}{n} \sin\frac{4\pi k}{n} \cdots \sin\frac{(n-1)\pi k}{n} = \pm \sqrt{n}.}
The sign in this formula, as in (15), still depends on $k$; but it is the positive one when $k = 1$. For the rest, we shall derive one more important theorem concerning this sign, which will be investigated more closely later on.
